\documentclass[12pt,a4paper]{report}

% ── Packages ──────────────────────────────────────────────────────────────────
\usepackage[T1]{fontenc}
\usepackage[utf8]{inputenc}
\usepackage[margin=2.5cm]{geometry}
\usepackage{lmodern}
\usepackage{microtype}
\usepackage{hyperref}
\usepackage{booktabs}
\usepackage{longtable}
\usepackage{array}
\usepackage{enumitem}
\usepackage{titlesec}
\usepackage{fancyhdr}
\usepackage{xcolor}
\usepackage{graphicx}
\usepackage{tcolorbox}
\usepackage{parskip}

% ── Colours ───────────────────────────────────────────────────────────────────
\definecolor{broadblue}{RGB}{0,60,120}
\definecolor{lightgray}{RGB}{245,245,245}
\definecolor{warnyellow}{RGB}{255,243,205}
\definecolor{warnborder}{RGB}{204,153,0}
\definecolor{noteblue}{RGB}{219,234,254}
\definecolor{noteborder}{RGB}{59,130,246}
\definecolor{codebg}{RGB}{245,245,245}

% ── Hyperref ──────────────────────────────────────────────────────────────────
\hypersetup{
  colorlinks=true,
  linkcolor=broadblue,
  urlcolor=broadblue,
  pdftitle={DashyDashboard Screenshot Attestation Setup Guide},
  pdfauthor={Broadridge},
  pdfsubject={Screenshot Attestation, Email and Storage Configuration}
}

% ── Headers / footers ─────────────────────────────────────────────────────────
\pagestyle{fancy}
\fancyhf{}
\fancyhead[L]{\small\color{broadblue}DashyDashboard — Screenshot Setup Guide}
\fancyhead[R]{\small\color{broadblue}Broadridge}
\fancyfoot[C]{\small\thepage}
\renewcommand{\headrulewidth}{0.4pt}

% ── Section formatting ────────────────────────────────────────────────────────
\titleformat{\chapter}[display]
  {\normalfont\Large\bfseries\color{broadblue}}
  {Chapter \thechapter}{10pt}{\Huge}
\titlespacing*{\chapter}{0pt}{20pt}{20pt}

\titleformat{\section}
  {\normalfont\large\bfseries\color{broadblue}}
  {\thesection}{1em}{}

% ── Custom boxes ──────────────────────────────────────────────────────────────
\tcbuselibrary{skins,breakable}

\newtcolorbox{notebox}{
  colback=noteblue,
  colframe=noteborder,
  fonttitle=\bfseries,
  title=Note,
  breakable
}

\newtcolorbox{warnbox}{
  colback=warnyellow,
  colframe=warnborder,
  fonttitle=\bfseries,
  title=Important,
  breakable
}

% A simple monospace code block on a light-gray background, for copy-pasteable commands.
\newtcolorbox{codebox}{
  colback=codebg,
  colframe=lightgray,
  boxrule=0.4pt,
  left=6pt,right=6pt,top=6pt,bottom=6pt,
  breakable
}

% ── Misc ──────────────────────────────────────────────────────────────────────
\setlist[enumerate]{itemsep=4pt,parsep=0pt}
\setlist[itemize]{itemsep=4pt,parsep=0pt}

% ═════════════════════════════════════════════════════════════════════════════
\begin{document}

% ── Title page ────────────────────────────────────────────────────────────────
\begin{titlepage}
  \centering
  \vspace*{3cm}
  {\color{broadblue}\rule{\linewidth}{1.5pt}}\\[0.6cm]
  {\Huge\bfseries\color{broadblue} DashyDashboard}\\[0.4cm]
  {\LARGE Screenshot Attestation — Setup \& Operations Guide}\\[0.3cm]
  {\color{broadblue}\rule{\linewidth}{1.5pt}}\\[1cm]
  {\large Screenshot storage, e-mail notifications, and database migrations}\\[2cm]
  {\large\itshape Broadridge Financial Solutions}\\[0.5cm]
  {\large June 2026}\\[3cm]
  \vfill
  {\small
    \begin{tabular}{ll}
      \textbf{Document type:}   & Setup / Operations Supplement \\
      \textbf{Companion to:}    & DashyDashboard Deployment Guide \\
      \textbf{Audience:}        & Server Administrators \& DBAs \\
      \textbf{Classification:}  & Internal \\
    \end{tabular}
  }
\end{titlepage}

% ── Table of contents ─────────────────────────────────────────────────────────
\tableofcontents
\clearpage

% ═════════════════════════════════════════════════════════════════════════════
\chapter{Overview}

This guide is a supplement to the main \emph{DashyDashboard Deployment Guide}. It covers the
operational steps required to enable the \textbf{screenshot attestation} feature and the associated
\textbf{e-mail notifications} introduced in the June 2026 release. Read the main deployment guide
first; this document only describes what is \emph{new}.

\section{What changed}

\begin{itemize}
  \item Associates now attach a screenshot to each tool they used during a cycle. Reviewers
        (managers, GFH/GFH-Delegate, and Admin) approve or reject each screenshot.
  \item Screenshots are written to a disk folder \emph{outside} the deployed application directory.
  \item Two new configuration sections --- \texttt{Screenshots} and \texttt{Email} --- must be
        present in \texttt{appsettings.Production.json} (or supplied via environment variables).
  \item Two new Entity Framework migrations add columns to the database.
  \item Optional event-driven e-mail (rejected screenshot / all approved) requires an internal SMTP
        relay and an approved sending mailbox.
\end{itemize}

\section{What an administrator must do (checklist)}

\begin{enumerate}
  \item Add the \texttt{Screenshots} and \texttt{Email} sections to
        \texttt{appsettings.Production.json}, or set the equivalent environment variables
        (Chapter~\ref{ch:config}).
  \item Create the screenshot storage folder outside the app folder and grant the IIS app-pool
        identity \emph{Modify} access (Chapter~\ref{ch:storage}).
  \item Apply the two database migrations, in order (Chapter~\ref{ch:migrations}).
  \item (Optional) Request the internal SMTP relay details and an approved sending mailbox, then
        enable e-mail (Chapter~\ref{ch:email}).
  \item Confirm the storage folder is covered by server backups and understand the retention policy
        (Chapter~\ref{ch:retention}).
\end{enumerate}

\begin{warnbox}
\texttt{dotnet publish} OVERWRITES \texttt{appsettings.Production.json} and the entire app folder on
every deployment. The screenshot folder must therefore live \emph{outside} the app folder, and the
\texttt{Screenshots}/\texttt{Email} configuration must be re-applied (or supplied via environment
variables, which survive a republish) after \emph{every} publish. See Chapter~\ref{ch:config}.
\end{warnbox}

% ═════════════════════════════════════════════════════════════════════════════
\chapter{Configuration Keys}
\label{ch:config}

The application reads two new configuration sections. The base \texttt{appsettings.json} ships with
safe defaults (storage on \texttt{C:}, e-mail disabled), but production should override at least
\texttt{Screenshots:RootPath} and, if mail is wanted, the \texttt{Email} section.

\section{The \texttt{Screenshots} section}

\begin{longtable}{>{\ttfamily\small}p{5.8cm} >{\raggedright\arraybackslash}p{8.2cm}}
\toprule
\normalfont\textbf{Key} & \textbf{Meaning} \\
\midrule
\endhead
Screenshots:RootPath &
  Absolute disk path under which all screenshots are stored, using the layout
  \texttt{\{root\}\textbackslash\{cycleId\}\textbackslash\{associateId\}\textbackslash}
  \texttt{\{clientId\}\textbackslash\{toolId\}.jpg}
  (plus a \texttt{\_thumb.jpg} thumbnail when JPEG is configured). \textbf{Must be OUTSIDE the deployed app folder}
  (Chapter~\ref{ch:storage}). The IIS app-pool identity needs \emph{Modify} permission here. \\
Screenshots:Format &
  Server-side storage format. Accepted values are \texttt{jpeg}, \texttt{jpg}, and \texttt{webp}.
  The production default is \texttt{jpeg}. Existing files remain readable after this value changes. \\
Screenshots:Quality &
  Main-image encoding quality from \texttt{1} to \texttt{100}. Production default \texttt{85}. \\
Screenshots:MaxLongEdge &
  Maximum width or height of the stored main image in pixels. Larger uploads are scaled down
  proportionally. Production default \texttt{1600}. Smaller images are not enlarged. \\
Screenshots:ThumbnailQuality &
  Thumbnail encoding quality from \texttt{1} to \texttt{100}. Production default \texttt{80}. \\
Screenshots:ThumbnailWidth &
  Maximum thumbnail width in pixels. Production default \texttt{200}. \\
Screenshots:RetainCycles &
  Integer number of recent cycles to keep on disk. Default \texttt{6}. Retention is a policy target
  only --- the automatic sweep is a future job (Chapter~\ref{ch:retention}). \\
\bottomrule
\end{longtable}

\section{The \texttt{Email} section}

\begin{longtable}{>{\ttfamily\small}p{6.0cm} >{\raggedright\arraybackslash}p{8.0cm}}
\toprule
\normalfont\textbf{Key} & \textbf{Meaning} \\
\midrule
\endhead
Email:Enabled &
  Master switch. \texttt{false} (the default) means no mail is ever sent --- the request paths run
  normally, mail is simply skipped. Set \texttt{true} only once the relay and sender are confirmed. \\
Email:SmtpHost &
  Hostname of the internal SMTP relay. The relay is contacted \emph{unauthenticated} and without TLS
  (corporate internal relay). \\
Email:SmtpPort &
  Relay port. Default \texttt{25}. \\
Email:From &
  The sending mailbox address that appears in the \emph{From} header. Must be an approved sender on
  the relay. \\
Email:AppUrl &
  Absolute base URL of the deployed app (e.g.\
  \texttt{http://clipvwbpod02/dashydashboard/}). Used to build the ``open DashyDashboard'' links in
  mail bodies. \\
Email:Events:ScreenshotRejected &
  Per-event toggle for the ``your screenshot was rejected'' mail to the associate. Default
  \texttt{true} (gated by the master switch). \\
Email:Events:AllApproved &
  Per-event toggle for the ``all your screenshots are approved'' mail. Sent exactly once, on the
  not-complete to complete transition. Default \texttt{true}. \\
\bottomrule
\end{longtable}

\section{Example \texttt{appsettings.Production.json} fragment}

Add these two sections alongside the existing \texttt{ConnectionStrings} block:

\begin{codebox}
\ttfamily\small
"Screenshots": \{\\
\hspace*{2em}"RootPath": "D:\textbackslash\textbackslash DashyData\textbackslash\textbackslash Screenshots",\\
\hspace*{2em}"Format": "jpeg",\\
\hspace*{2em}"Quality": 85,\\
\hspace*{2em}"MaxLongEdge": 1600,\\
\hspace*{2em}"ThumbnailQuality": 80,\\
\hspace*{2em}"ThumbnailWidth": 200,\\
\hspace*{2em}"RetainCycles": 6\\
\},\\
"Email": \{\\
\hspace*{2em}"Enabled": false,\\
\hspace*{2em}"SmtpHost": "smtprelay.internal.example.com",\\
\hspace*{2em}"SmtpPort": 25,\\
\hspace*{2em}"From": "dashydashboard-noreply@example.com",\\
\hspace*{2em}"AppUrl": "http://clipvwbpod02/dashydashboard/",\\
\hspace*{2em}"Events": \{ "ScreenshotRejected": true, "AllApproved": true \}\\
\}
\end{codebox}

\begin{warnbox}
Because \texttt{dotnet publish} overwrites \texttt{appsettings.Production.json}, these two sections
must be re-added after \emph{every} publish --- this extends the existing
``Restore \texttt{appsettings.Production.json}'' step in the deployment checklist. The most reliable
way to avoid re-editing the file each time is to use environment variables instead (next section);
they live outside the published artifact and survive a republish.
\end{warnbox}

\section{Environment-variable alternative (recommended for per-machine values)}

Every configuration key can be supplied as a machine environment variable instead of being written
into the JSON file. ASP.NET Core maps a nested key to an environment variable by replacing each
colon (\texttt{:}) with a double underscore (\texttt{\_\_}):

\begin{longtable}{>{\ttfamily}p{6.6cm} >{\ttfamily}p{7.4cm}}
\toprule
\normalfont\textbf{Configuration key} & \normalfont\textbf{Environment variable} \\
\midrule
\endhead
Screenshots:RootPath          & Screenshots\_\_RootPath \\
Screenshots:Format            & Screenshots\_\_Format \\
Screenshots:Quality           & Screenshots\_\_Quality \\
Screenshots:MaxLongEdge       & Screenshots\_\_MaxLongEdge \\
Screenshots:ThumbnailQuality  & Screenshots\_\_ThumbnailQuality \\
Screenshots:ThumbnailWidth    & Screenshots\_\_ThumbnailWidth \\
Screenshots:RetainCycles      & Screenshots\_\_RetainCycles \\
Email:Enabled                 & Email\_\_Enabled \\
Email:SmtpHost                & Email\_\_SmtpHost \\
Email:From                    & Email\_\_From \\
Email:Events:ScreenshotRejected & Email\_\_Events\_\_ScreenshotRejected \\
\bottomrule
\end{longtable}

\begin{notebox}
Prefer the environment-variable form for \emph{per-machine} values --- above all
\texttt{Screenshots\_\_RootPath}. It mirrors the existing \texttt{ConnectionStrings\_\_Default}
pattern already used on these servers: the value is set once at the machine level, lives outside the
deployed artifact, and therefore survives every \texttt{dotnet publish} without any file editing. An
environment variable takes precedence over the same key in \texttt{appsettings*.json}.
\end{notebox}

% ═════════════════════════════════════════════════════════════════════════════
\chapter{Screenshot Storage Folder}
\label{ch:storage}

\section{Why it must live outside the app folder}

Each deployment replaces the entire contents of the application folder
(\texttt{C:\textbackslash inetpub\textbackslash wwwroot\textbackslash DashyDashboard}). If
screenshots were stored inside that folder they would be deleted on the next deploy. The storage
root must therefore be a separate location, for example:

\begin{codebox}
\ttfamily D:\textbackslash DashyData\textbackslash Screenshots
\end{codebox}

\section{One-time elevated setup (per server)}

This is a one-time step on each server, performed by an administrator with elevation.

\begin{enumerate}
  \item Create the folder (substitute your chosen path):
  \begin{codebox}
  \ttfamily mkdir D:\textbackslash DashyData\textbackslash Screenshots
  \end{codebox}

  \item Grant the IIS application-pool identity \emph{Modify} permission on the folder, applied to
        the folder, its subfolders, and files. Use \texttt{icacls} from an elevated Command Prompt
        or PowerShell. For the local development pool (\texttt{DashyDashboardPool}):
  \begin{codebox}
  \ttfamily icacls "D:\textbackslash DashyData\textbackslash Screenshots" /grant "IIS AppPool\textbackslash DashyDashboardPool:(OI)(CI)M"
  \end{codebox}

  \item On the production server, substitute the production application pool's name for
        \texttt{DashyDashboardPool}. Confirm the exact pool name in IIS Manager
        (Application Pools) before running the command.
\end{enumerate}

\begin{notebox}
The \texttt{icacls} flags mean: \texttt{(OI)} object inherit (files in the folder),
\texttt{(CI)} container inherit (subfolders), and \texttt{M} the Modify permission set
(read, write, create, delete). The app needs to create per-cycle subfolders and write image files,
so Modify --- not just Write --- is required.
\end{notebox}

\begin{warnbox}
The login name passed to \texttt{icacls} is derived from the application-pool name and is always
prefixed \texttt{IIS AppPool\textbackslash}. It is \emph{not} a domain account and will not appear
when browsing for users --- type it exactly. If the application logs ``Access to the path
\ldots\ is denied'' when an associate uploads a screenshot, this ACL grant is the first thing to
check.
\end{warnbox}

\section{Setting the path}

After the folder exists, point the application at it via either:

\begin{itemize}
  \item \texttt{Screenshots:RootPath} in \texttt{appsettings.Production.json}, \emph{or}
  \item the machine environment variable \texttt{Screenshots\_\_RootPath} (recommended ---
        Chapter~\ref{ch:config}). The app will refuse to start if no root path is configured.
\end{itemize}

% ═════════════════════════════════════════════════════════════════════════════
\chapter{Database Migrations}
\label{ch:migrations}

Two Entity Framework migrations must be applied, \textbf{in this order} (this is their merge order ---
the first shipped with the Tool-User-ID feature, the second with the screenshot feature):

\begin{enumerate}
  \item \texttt{20260611112843\_AddToolUserIdToUsersToolAccess} --- adds
        \texttt{UsersToolAccess.ToolUserId} (\texttt{nvarchar(100) NULL}).
  \item \texttt{20260611115259\_AddScreenshotColumnsToToolCycleAttestation} --- adds seven nullable
        columns to \texttt{ToolCycleAttestation}: \texttt{ScreenshotPath}, \texttt{ScreenshotHash},
        \texttt{ScreenshotUploadedAt}, \texttt{ScreenshotStatus}, \texttt{ScreenshotReviewedBy},
        \texttt{ScreenshotReviewedAt}, and \texttt{ScreenshotRejectReason}.
\end{enumerate}

\begin{warnbox}
Apply both migrations in the order listed. Applying the screenshot migration before the
Tool-User-ID migration will leave the database history out of step with the deployed binaries.
\end{warnbox}

\section{Option A --- apply directly with the EF tools (dev / test)}

From the API project directory, with the connection string pointing at the target database:

\begin{codebox}
\ttfamily\small
\$env:ConnectionStrings\_\_Default = "Server=\ldots;Database=\ldots;\ldots"\\
dotnet ef database update
\end{codebox}

\texttt{dotnet ef database update} (with no migration name) brings the database fully up to date,
applying any pending migrations in order. This is appropriate for the development and test databases.

\section{Option B --- generate an idempotent SQL script for the DBA (production)}

On production, the database is usually changed only by a DBA running a reviewed script. Generate an
idempotent script that can be run safely regardless of which migrations are already applied:

\begin{codebox}
\ttfamily\small
dotnet ef migrations script --idempotent --output screenshots-migrations.sql
\end{codebox}

Hand \texttt{screenshots-migrations.sql} to the DBA to execute against \texttt{ClientsAppAttestation}
in SSMS. An idempotent script wraps each migration in a guard so already-applied migrations are
skipped --- it is safe to run even if the first migration was applied earlier.

\begin{notebox}
The application's app-pool SQL login has \texttt{db\_ddladmin}, so EF can also self-apply migrations
on first launch in environments where that is the chosen workflow. On production, prefer the
reviewed idempotent script (Option B).
\end{notebox}

% ═════════════════════════════════════════════════════════════════════════════
\chapter{E-mail Notifications}
\label{ch:email}

E-mail is \textbf{disabled by default} (\texttt{Email:Enabled = false}). When disabled, the
application behaves identically except that no mail is sent --- there is no error and no impact on
the request paths. Enable it only after the relay and sending mailbox are confirmed.

\section{Events}

\begin{itemize}
  \item \textbf{ScreenshotRejected} --- sent to the associate when a reviewer rejects one of their
        screenshots. Includes the cycle, client, tool, reviewer name, the rejection reason, and a
        link to re-upload in DashyDashboard.
  \item \textbf{AllApproved} --- sent to the associate exactly once, when the last of their
        screenshots is approved and their attestation becomes complete.
\end{itemize}

Delivery is fire-and-forget: events are queued in-process and sent by a background consumer over an
unauthenticated internal SMTP relay. A misconfigured or unreachable relay never blocks or fails a
user request.

\section{IT request template}

The relay host/port and an approved sending mailbox are not Broadridge-developer-owned --- request
them from the messaging/IT team. A short copy-pasteable request:

\begin{codebox}
\ttfamily\small
Subject: SMTP relay + sending mailbox for DashyDashboard\\[4pt]
We are deploying an internal web application (DashyDashboard) on server\\
CLIPVWBPOD02 that needs to send automated, plain-text notification e-mails\\
to associates (screenshot rejected / attestation complete).\\[4pt]
Please provide:\\
\hspace*{1em}1. The internal SMTP relay hostname and port that this server may\\
\hspace*{3em}use to send mail (unauthenticated relay from this host is fine).\\
\hspace*{1em}2. An approved "from" sending mailbox / display address for the\\
\hspace*{3em}From header (e.g. dashydashboard-noreply@<domain>).\\[4pt]
Volume is low (a handful of mails per attestation cycle). No inbound mail\\
or mailbox monitoring is required.
\end{codebox}

\section{Enabling}

Once you have the relay details and sending mailbox:

\begin{enumerate}
  \item Set \texttt{Email:SmtpHost}, \texttt{Email:SmtpPort}, and \texttt{Email:From} (and confirm
        \texttt{Email:AppUrl} points at the live site).
  \item Set \texttt{Email:Enabled} to \texttt{true}.
  \item Leave the per-event toggles at \texttt{true} unless a specific event is to be suppressed.
\end{enumerate}

% ═════════════════════════════════════════════════════════════════════════════
\chapter{Retention \& Backup}
\label{ch:retention}

\section{Retention policy}

Cycles are monthly. The policy is to retain the most recent \textbf{6} cycles' screenshots on disk
(\texttt{Screenshots:RetainCycles = 6}, roughly six months).

\begin{warnbox}
The automatic retention sweep is a \emph{future} scheduled job and is \textbf{not yet running}.
Until it ships, retention is manual.
\end{warnbox}

\section{Manual purge (until the sweep ships)}

Each cycle's screenshots live under a top-level folder named for the numeric cycle ID:

\begin{codebox}
\ttfamily \{root\}\textbackslash\{cycleId\}\textbackslash\ldots
\end{codebox}

To purge, delete the \texttt{\{root\}\textbackslash\{cycleId\}} folders for cycles older than the
most recent six. Identify cycle IDs from \texttt{dbo.Cycles}. Deleting a folder removes only the
images; the attestation rows and their stored paths remain in the database.

\section{Backup coverage}

The screenshot root folder holds attestation evidence that exists nowhere else. It should be added
to the server's backup coverage so that an evidence loss is recoverable. Whether to back it up, and
the retention of those backups, is an ops decision --- this guide simply flags that the folder is
outside the app folder and outside the database, so it is easy to miss.

% ═════════════════════════════════════════════════════════════════════════════
\chapter{Local Development Box (\texttt{deploy-dev.ps1})}
\label{ch:devbox}

The local PRV-SERVER box runs the production-mirror artifact under IIS against a local SQL instance.
For the screenshot feature to work there, the same two prerequisites apply as on any server: a
storage folder outside the app folder, and the configuration pointing at it. These should be added
to the local \texttt{deploy-dev.ps1} (the elevated, machine-specific deploy script). The changes are
described here; perform them in an elevated session --- they are not part of the unelevated build.

\begin{enumerate}
  \item \textbf{Set the machine environment variable} so it survives every republish, mirroring the
        existing \texttt{ConnectionStrings\_\_Default} line:
  \begin{codebox}
  \ttfamily\small
  [Environment]::SetEnvironmentVariable(\\
  \hspace*{2em}"Screenshots\_\_RootPath", "D:\textbackslash DashyData\textbackslash Screenshots", "Machine")
  \end{codebox}

  \item \textbf{Create the folder and grant the local app-pool identity Modify} (idempotent):
  \begin{codebox}
  \ttfamily\small
  New-Item -ItemType Directory -Force "D:\textbackslash DashyData\textbackslash Screenshots"\\
  icacls "D:\textbackslash DashyData\textbackslash Screenshots" /grant "IIS AppPool\textbackslash DashyDashboardPool:(OI)(CI)M"
  \end{codebox}
\end{enumerate}

\begin{notebox}
\texttt{deploy-dev.ps1} already performs the analogous elevated work for the connection string
(setting \texttt{ConnectionStrings\_\_Default} as a machine variable) and for IIS. Adding the two
steps above keeps the local box self-contained and republish-proof. The local pool is
\texttt{DashyDashboardPool}; on production substitute the real pool name.
\end{notebox}

% ═════════════════════════════════════════════════════════════════════════════
\chapter{Future Phase}
\label{ch:future}

The following are documented as the planned \emph{next} phase. They reuse the e-mail plumbing
(the queue + send split in \texttt{EmailService}) and the storage service's \texttt{Delete} method
that already exist, so no new infrastructure is required --- only a scheduler.

\begin{itemize}
  \item \textbf{Scheduled \texttt{BackgroundService}.} A single hosted service that wakes on a timer
        (e.g.\ daily) and performs the periodic tasks below. It composes mail via the existing
        \texttt{EmailService} and enqueues it the same way the request paths do.
  \item \textbf{Due-date reminder mails.} Seven days before a cycle's \texttt{DueDate}, mail
        associates who have not yet completed their attestation.
  \item \textbf{Manager pending-approvals digest.} A periodic summary to each reviewer listing the
        screenshots awaiting their approval.
  \item \textbf{Retention sweep.} Automatically delete \texttt{\{root\}\textbackslash\{cycleId\}}
        folders for cycles older than \texttt{Screenshots:RetainCycles}, removing the manual purge
        step in Chapter~\ref{ch:retention}.
\end{itemize}

\begin{notebox}
None of these are built in the current release. They are listed so that the configuration keys
(\texttt{Screenshots:RetainCycles}, the \texttt{Email} section) and the storage layout described in
this guide are understood as forward-compatible with the next phase.
\end{notebox}

% ═════════════════════════════════════════════════════════════════════════════
\end{document}
